\chapter{Conclusion}
\label{ch:conclusion}

This thesis presented the first implementation and empirical study of
the Chen--Wein--Zhang next-to-shortest-path algorithm. The algorithm
resolves a thirty-year-old open question for directed graphs with
strictly positive edge weights and stands out as the only NSP algorithm
with a fixed polynomial bound on general inputs.

\section{Summary of findings}

\begin{enumerate}
  \item The algorithm is implementable from the paper's pseudocode
    without needing access to optimised 2-VDP subroutines: a
    straightforward max-flow-based 2-VDP yields a working implementation
    at the cost of a higher constant factor (the asymptotic
    $O(|V|+|E|)$ bound is the same as Tholey's linear-time algorithm).
  \item On the four benign random graph families we measured
    (Erd\H{o}s--R\'enyi, random DAG, layered, grid) at $|V|\le 50$,
    Yen's $k$-shortest-simple-paths algorithm is several orders of
    magnitude faster than CWZ. This was expected: our random generators
    almost never produce many distinct simple shortest paths, so Yen
    rarely takes more than a handful of iterations.
  \item On a deliberate adversarial family of \emph{diamond-chain}
    graphs designed to have $2^{k}$ distinct shortest $s\to t$ paths
    (Section~\ref{sec:adversarial}), the picture reverses dramatically.
    At $k=15$ ($|V|=46$, $|E|=61$) Yen takes $59$~seconds while CWZ
    finishes in $1.4$~milliseconds --- a ${\sim}41{,}000\times$ ratio
    that grows exponentially with $k$. CWZ scales smoothly to $k=24$
    ($|V|=73$) in tens of microseconds; Yen would require many hours on
    the same instance by extrapolation. Worse, with a practical
    enumeration cap Yen does not merely slow down: at $|V|\ge 40$ it
    exceeds its cap and returns the \emph{wrong} answer (``no NSP'')
    while CWZ stays correct. This is the empirical demonstration of why
    a polynomial-time NSP algorithm matters.
  \item Our implementation is a faithful transcription of the paper's
    three-stage algorithm: \textsc{NextSP} (fold off-shortest-path
    vertices), \textsc{NextSP-Straight} (remove forward/sideways
    back-edges, recording their candidate paths, to obtain a strictly
    $(s,t)$-layered graph), and \textsc{NextSP-Layered} (the Lemma~5.3
    $6$-tuple enumeration). It matches the brute-force oracle with
    \emph{zero} mismatches across more than $35{,}000$ random instances
    ($|V|\le 14$) and matches Yen on every benign benchmark instance up
    to $|V|=50$. An ablation confirms the $6$-tuple enumeration alone
    suffices on the strictly-layered graph --- the property the paper's
    correctness proof relies on. The one remaining approximation is the
    $K=5$ multi-edge cap in \textsc{NextSP} (below).
\end{enumerate}

\section{Limitations and future work}

\paragraph{Engineering improvements.} Replacing the max-flow 2-VDP by
Tholey's linear-time algorithm~\cite{Tholey2012} would lower the
constant factor on the inner subroutine. Eliminating the
\textsc{NextSP-Straight} subdivision via clever distance indexing,
suggested by the paper's authors as a clarity-oriented choice, could
provide further speedup.

\paragraph{Larger graphs.} Our experiments stop at $|V|=50$. The
empirical curve on Erd\H{o}s--R\'enyi suggests the actual runtime grows
polynomially with an exponent of roughly $3$, well below the
worst-case $|V|^{4}|E|^{3}$ bound; characterising this gap on still
larger or differently-structured inputs would be valuable.

\paragraph{Multi-edge cap tightening.} The $K=5$ multi-edge cap in
\texttt{reduce\_to\_straight} is the one approximation in our pipeline.
The faithful fix is to implement \textsc{NextSP}'s own
candidate-recording step (mirroring what \textsc{NextSP-Straight}
already does for back-edges in our code): record the candidate path
through each folded vertex, then keep only the single minimum-weight
shortcut per pair. This would eliminate both the cap and the Cartesian
multi-edge blow-up that motivated it, making the entire pipeline a
proven transcription of the paper.

\paragraph{Adversarial benchmarks.} The diamond-chain family is the
simplest construction we know of that defeats Yen. Richer adversarial
graphs --- recursive nestings, irregular branching factors, weight
ratios that vary with depth --- would map out the boundary between
``Yen survives'' and ``CWZ is the only viable option'' more finely.

\paragraph{Algorithmic refinements.} The paper's authors note their
bound is ``not optimised''. Candidate optimisations include exploiting
layer structure to prune the $6$-tuple enumeration and amortising 2-VDP
queries across nearby tuples, as well as replacing our max-flow 2-VDP
with Tholey's linear-time routine.
